\documentclass[10pt,a4paper]{amsart}
\usepackage{amsmath,amssymb,amsthm,graphicx}
% Retained minimal preamble for the Stacks Project chapter flat.tex (More on Flatness).
% Verbatim excerpt of the chapter preamble: Xy-pic setup (native preamble lines 14-18) and the native
% theorem declarations with the mathematical macros (native preamble lines 154-217). Copied mechanically;
% nothing added, removed or redefined. The wrapper class and the amsmath/amssymb/amsthm/graphicx base
% packages are supplied by the assembler, so the class/externaldocument part of the native preamble is omitted.
% For commutative diagrams we use Xy-pic
\usepackage[all]{xy}

% We use 2cell for 2-commutative diagrams.
\xyoption{2cell}
% Theorem environments.
%
\theoremstyle{plain}
\newtheorem{theorem}[subsection]{Theorem}
\newtheorem{proposition}[subsection]{Proposition}
\newtheorem{lemma}[subsection]{Lemma}

\theoremstyle{definition}
\newtheorem{definition}[subsection]{Definition}
\newtheorem{example}[subsection]{Example}
\newtheorem{exercise}[subsection]{Exercise}
\newtheorem{situation}[subsection]{Situation}

\theoremstyle{remark}
\newtheorem{remark}[subsection]{Remark}
\newtheorem{remarks}[subsection]{Remarks}

\numberwithin{equation}{subsection}

% Macros
%
\def\lim{\mathop{\mathrm{lim}}\nolimits}
\def\colim{\mathop{\mathrm{colim}}\nolimits}
\def\Spec{\mathop{\mathrm{Spec}}}
\def\Hom{\mathop{\mathrm{Hom}}\nolimits}
\def\Ext{\mathop{\mathrm{Ext}}\nolimits}
\def\SheafHom{\mathop{\mathcal{H}\!\mathit{om}}\nolimits}
\def\SheafExt{\mathop{\mathcal{E}\!\mathit{xt}}\nolimits}
\def\Sch{\mathit{Sch}}
\def\Mor{\mathop{\mathrm{Mor}}\nolimits}
\def\Ob{\mathop{\mathrm{Ob}}\nolimits}
\def\Sh{\mathop{\mathit{Sh}}\nolimits}
\def\NL{\mathop{N\!L}\nolimits}
\def\CH{\mathop{\mathrm{CH}}\nolimits}
\def\proetale{{pro\text{-}\acute{e}tale}}
\def\etale{{\acute{e}tale}}
\def\QCoh{\mathit{QCoh}}
\def\Ker{\mathop{\mathrm{Ker}}}
\def\Im{\mathop{\mathrm{Im}}}
\def\Coker{\mathop{\mathrm{Coker}}}
\def\Coim{\mathop{\mathrm{Coim}}}

% Boxtimes
%
\DeclareMathSymbol{\boxtimes}{\mathbin}{AMSa}{"02}

%
% Macros for moduli stacks/spaces
%
\def\QCohstack{\mathcal{QC}\!\mathit{oh}}
\def\Cohstack{\mathcal{C}\!\mathit{oh}}
\def\Spacesstack{\mathcal{S}\!\mathit{paces}}
\def\Quotfunctor{\mathrm{Quot}}
\def\Hilbfunctor{\mathrm{Hilb}}
\def\Curvesstack{\mathcal{C}\!\mathit{urves}}
\def\Polarizedstack{\mathcal{P}\!\mathit{olarized}}
\def\Complexesstack{\mathcal{C}\!\mathit{omplexes}}
% \Pic is the operator that assigns to X its picard group, usage \Pic(X)
% \Picardstack_{X/B} denotes the Picard stack of X over B
% \Picardfunctor_{X/B} denotes the Picard functor of X over B
\def\Pic{\mathop{\mathrm{Pic}}\nolimits}
\def\Picardstack{\mathcal{P}\!\mathit{ic}}
\def\Picardfunctor{\mathrm{Pic}}
\def\Deformationcategory{\mathcal{D}\!\mathit{ef}}
\usepackage{cleveref}
\crefname{originaldefinition}{original definition}{original definitions}
\Crefname{originaldefinition}{Original definition}{Original definitions}
\makeatletter
\newlabel{situation-flat-dimension-n}{{05MT}{source}{}{sourcepointer-1}{}}
\newlabel{situation-flat-dimension-n@cref}{{[originaldefinition][0][]05MT}{[1][0][]source}{}{}{}}
\makeatother
\crefname{originalequation}{original equation}{original equations}
\Crefname{originalequation}{Original equation}{Original equations}
\makeatletter
\newlabel{equation-flat-dimension-n}{{05MU}{source}{}{sourcepointer-2}{}}
\newlabel{equation-flat-dimension-n@cref}{{[originalequation][0][]05MU}{[1][0][]source}{}{}{}}
\makeatother
\crefname{originaldefinition}{original definition}{original definitions}
\Crefname{originaldefinition}{Original definition}{Original definitions}
\makeatletter
\newlabel{definition-complete-devissage}{{05HH}{source}{}{sourcepointer-3}{}}
\newlabel{definition-complete-devissage@cref}{{[originaldefinition][0][]05HH}{[1][0][]source}{}{}{}}
\makeatother
\crefname{originaldefinition}{original definition}{original definitions}
\Crefname{originaldefinition}{Original definition}{Original definitions}
\makeatletter
\newlabel{definition-one-step-devissage}{{05H4}{source}{}{sourcepointer-4}{}}
\newlabel{definition-one-step-devissage@cref}{{[originaldefinition][0][]05H4}{[1][0][]source}{}{}{}}
\makeatother
\crefname{originallemma}{original lemma}{original lemmas}
\Crefname{originallemma}{Original lemma}{Original lemmas}
\makeatletter
\newlabel{lemma-induction-step-fp}{{05IA}{source}{}{sourcepointer-5}{}}
\newlabel{lemma-induction-step-fp@cref}{{[originallemma][0][]05IA}{[1][0][]source}{}{}{}}
\makeatother
\crefname{originallemma}{original lemma}{original lemmas}
\Crefname{originallemma}{Original lemma}{Original lemmas}
\makeatletter
\newlabel{lemma-compare-H-F}{{05UE}{source}{}{sourcepointer-6}{}}
\newlabel{lemma-compare-H-F@cref}{{[originallemma][0][]05UE}{[1][0][]source}{}{}{}}
\makeatother
\crefname{originallemma}{original lemma}{original lemmas}
\Crefname{originallemma}{Original lemma}{Original lemmas}
\makeatletter
\newlabel{lemma-free-at-generic-points-representable}{{05UC}{source}{}{sourcepointer-7}{}}
\newlabel{lemma-free-at-generic-points-representable@cref}{{[originallemma][0][]05UC}{[1][0][]source}{}{}{}}
\makeatother
\title{Representability of the functor of points where a finitely presented module is flat in dimensions at least d, given a complete devissage}
\author{The Stacks Project Authors}
\date{stacks-project:flat.tex:local-snapshot-20261011; GNU FDL 1.2 or later}
\begin{document}
\maketitle
\noindent\textit{Editorial note.} All delivered mathematical text is the author's own, copied from the pinned Stacks Project chapter file flat.tex; no statement or proof step was rewritten, reordered or removed.\par\medskip
\noindent\textit{Editorial note.} The preamble is the chapter's own retained minimal preamble (sources/stacks-flat/retained-preamble.tex), a verbatim excerpt of the native Xy-pic setup and of the native theorem declarations and mathematical macros; no macro was invented or redefined, and the wrapper class and the amsmath/amssymb/amsthm/graphicx base packages are supplied by the assembler.\par\medskip
\noindent\textit{Editorial note.} The counted statement and proof contain no cross reference, no citation and no included file, so the delivered body is self contained.\par\medskip
\noindent\textit{Editorial note.} Canonical location of the counted proposition: Stacks tag 05UF (sources/stacks-flat/pinned-original-source.tex line 7806).\par\medskip
\noindent\textit{Allowed standard prerequisites (editorial).}\begin{itemize}
\item The flatness-in-dimension-n situation, its equation, and the one-step and complete devissage definitions. Same chapter, flat.tex, Situation labelled situation-flat-dimension-n (source line 6311), Equation labelled equation-flat-dimension-n (source line 6325), Definition labelled definition-one-step-devissage (source line 780) and Definition labelled definition-complete-devissage (source line 1195), all invoked in the statement and at proof lines 7830-7836. Fix the objects, the functor F\_d and the devissage data (Z\_k, Y\_k, i\_k, pi\_k, G\_k, alpha\_k) with the smoothness of Y\_k over S and the closedness of i\_k.
\item The induction step making alpha\_k universally injective after shrinking. Same chapter, Lemma labelled lemma-induction-step-fp (source line 3141), invoked at proof lines 7867-7869. Supplies that after replacing S by an affine neighbourhood of s the map alpha\_1 is S-universally injective, which produces the displayed short exact sequence.
\item The boundary case d\_1 = d and the comparison of the two functors. Same chapter, Lemmas labelled lemma-compare-H-F (source line 7778) and lemma-free-at-generic-points-representable (source line 7631), invoked at proof lines 7888-7891. Supply the representability conclusion in the remaining boundary case; the induction and the functorial reductions are inside the counted span.
\end{itemize}\medskip
\ifdefined\hypertarget\hypertarget{sourcepointer-1}{}\fi
\noindent\textit{Source reference.} Original definition 05MT: sources/stacks-flat/pinned-original-source.tex, source lines 6311--6311. Same-chapter pointer into the author environment situation at source line 6311; invoked at proof line None of the counted proof. It is retained only as a navigation pointer; every load-bearing step is inside the delivered span.\par\medskip
\ifdefined\hypertarget\hypertarget{sourcepointer-2}{}\fi
\noindent\textit{Source reference.} Original equation 05MU: sources/stacks-flat/pinned-original-source.tex, source lines 6325--6325. Same-chapter pointer into the author environment equation at source line 6325; invoked at proof line None of the counted proof. It is retained only as a navigation pointer; every load-bearing step is inside the delivered span.\par\medskip
\ifdefined\hypertarget\hypertarget{sourcepointer-3}{}\fi
\noindent\textit{Source reference.} Original definition 05HH: sources/stacks-flat/pinned-original-source.tex, source lines 1195--1195. Same-chapter pointer into the author environment definition at source line 1195; invoked at proof line 7832 of the counted proof. It is retained only as a navigation pointer; every load-bearing step is inside the delivered span.\par\medskip
\ifdefined\hypertarget\hypertarget{sourcepointer-4}{}\fi
\noindent\textit{Source reference.} Original definition 05H4: sources/stacks-flat/pinned-original-source.tex, source lines 780--780. Same-chapter pointer into the author environment definition at source line 780; invoked at proof line 7835 of the counted proof. It is retained only as a navigation pointer; every load-bearing step is inside the delivered span.\par\medskip
\ifdefined\hypertarget\hypertarget{sourcepointer-5}{}\fi
\noindent\textit{Source reference.} Original lemma 05IA: sources/stacks-flat/pinned-original-source.tex, source lines 3141--3141. Same-chapter pointer into the author environment lemma at source line 3141; invoked at proof line 7868 of the counted proof. It is retained only as a navigation pointer; every load-bearing step is inside the delivered span.\par\medskip
\ifdefined\hypertarget\hypertarget{sourcepointer-6}{}\fi
\noindent\textit{Source reference.} Original lemma 05UE: sources/stacks-flat/pinned-original-source.tex, source lines 7778--7778. Same-chapter pointer into the author environment lemma at source line 7778; invoked at proof line 7889 of the counted proof. It is retained only as a navigation pointer; every load-bearing step is inside the delivered span.\par\medskip
\ifdefined\hypertarget\hypertarget{sourcepointer-7}{}\fi
\noindent\textit{Source reference.} Original lemma 05UC: sources/stacks-flat/pinned-original-source.tex, source lines 7631--7631. Same-chapter pointer into the author environment lemma at source line 7631; invoked at proof line 7891 of the counted proof. It is retained only as a navigation pointer; every load-bearing step is inside the delivered span.\par\medskip
\begin{lemma}
\label{lemma-flat-dimension-n-representable}
In Situation \ref{situation-flat-dimension-n}.
Let $s \in S$ let $d \geq 0$. Assume
\begin{enumerate}
\item there exists a complete d\'evissage
of $\mathcal{F}/X/S$ over some point $s \in S$,
\item $X$ is of finite presentation over $S$,
\item $\mathcal{F}$ is an $\mathcal{O}_X$-module of finite presentation, and
\item $\mathcal{F}$ is flat in dimensions $\geq d + 1$ over $S$.
\end{enumerate}
Then after possibly replacing $S$ by an open neighbourhood
of $s$ the functor $F_d$ (\ref{equation-flat-dimension-n})
is representable by a monomorphism $Z_d \to S$ of finite presentation.
\end{lemma}

\begin{proof}
A preliminary remark is that $X$, $S$ are affine schemes and that it
suffices to prove $F_d$ is representable by a monomorphism of finite
presentation $Z_d \to S$ on the category of affine schemes over $S$.
(Of course we do not require $Z_d$ to be affine.)
Hence throughout the proof of
the lemma we work in the category of affine schemes over $S$.

\medskip\noindent
Let $(Z_k, Y_k, i_k, \pi_k, \mathcal{G}_k, \alpha_k)_{k = 1, \ldots, n}$
be a complete d\'evissage of $\mathcal{F}/X/S$ over $s$, see
Definition \ref{definition-complete-devissage}.
We will use induction on the length $n$ of the d\'evissage.
Recall that $Y_k \to S$ is smooth with geometrically irreducible fibres, see
Definition \ref{definition-one-step-devissage}.
Let $d_k$ be the relative dimension of $Y_k$ over $S$.
Recall that $i_{k, *}\mathcal{G}_k = \Coker(\alpha_k)$ and
that $i_k$ is a closed immersion.
By the definitions referenced above we have
$d_1 = \dim(\text{Supp}(\mathcal{F}_s))$ and
$$
d_k = \dim(\text{Supp}(\Coker(\alpha_{k - 1})_s))
= \dim(\text{Supp}(\mathcal{G}_{k, s}))
$$
for $k = 2, \ldots, n$. It follows that $d_1 > d_2 > \ldots > d_n \geq 0$
because $\alpha_k$ is an isomorphism in the generic point of $(Y_k)_s$.

\medskip\noindent
Note that $i_1$ is a closed immersion and
$\mathcal{F} = i_{1, *}\mathcal{G}_1$.
Hence for any morphism of schemes $T \to S$ with $T$ affine,
we have $\mathcal{F}_T = i_{1, T, *}\mathcal{G}_{1, T}$ and
$i_{1, T}$ is still a closed immersion of schemes over $T$.
Thus $\mathcal{F}_T$ is flat in dimensions $\geq d$ over $T$
if and only if $\mathcal{G}_{1, T}$ is flat in dimensions $\geq d$ over $T$.
Because $\pi_1 : Z_1 \to Y_1$ is finite we see in the same manner that
$\mathcal{G}_{1, T}$ is flat in dimensions $\geq d$ over $T$
if and only if $\pi_{1, T, *}\mathcal{G}_{1, T}$ is flat in dimensions
$\geq d$ over $T$. The same arguments work for
``flat in dimensions $\geq d + 1$'' and we conclude in particular that
$\pi_{1, *}\mathcal{G}_1$ is flat over $S$ in dimensions $\geq d + 1$
by our assumption on $\mathcal{F}$.

\medskip\noindent
Suppose that $d_1 > d$. It follows from the discussion above that
in particular $\pi_{1, *}\mathcal{G}_1$ is flat over $S$ at
the generic point of $(Y_1)_s$. By
Lemma \ref{lemma-induction-step-fp}
we may replace $S$ by an affine neighbourhood of $s$ and assume that
$\alpha_1$ is $S$-universally injective. Because $\alpha_1$ is
$S$-universally injective, for any morphism $T \to S$ with $T$ affine,
we have a short exact sequence
$$
0 \to \mathcal{O}_{Y_{1, T}}^{\oplus r_1}
\to \pi_{1, T, *}\mathcal{G}_{1, T} \to \Coker(\alpha_1)_T \to 0
$$
and still the first arrow is $T$-universally injective. Hence the set
of points of $(Y_1)_T$ where $\pi_{1, T, *}\mathcal{G}_{1, T}$ is flat over
$T$ is the same as the set of points of $(Y_1)_T$ where
$\Coker(\alpha_1)_T$ is flat over $S$. In this way the question
reduces to the sheaf $\Coker(\alpha_1)$ which has a complete
d\'evissage of length $n - 1$ and we win by induction.

\medskip\noindent
If $d_1 < d$ then $F_d$ is represented by $S$ and we win.

\medskip\noindent
The last case is the case $d_1 = d$. This case follows from a combination of
Lemma \ref{lemma-compare-H-F}
and
Lemma \ref{lemma-free-at-generic-points-representable}.
\end{proof}

\end{document}
